\documentclass[aps,prb,onecolumn,nofootinbib,superscriptaddress]{revtex4-2}
\usepackage{amsmath,amssymb,amsthm,mathtools,bm}
\usepackage{booktabs}
\usepackage{tabularx}
\usepackage{xcolor}
\usepackage[colorlinks=true,linkcolor=blue,citecolor=blue,urlcolor=blue]{hyperref}
\usepackage{microtype}

\hypersetup{
  pdftitle={Open-System / Monitored-System Response Theory for a Class-A Adaptive Fermion Circuit},
  pdfauthor={Class-A fermionic adaptive-circuit project}
}
\newcolumntype{L}[1]{>{\raggedright\arraybackslash}p{#1}}
\newcolumntype{Y}{>{\raggedright\arraybackslash}X}

\newcommand{\Id}{\mathbf{1}}
\newcommand{\Tr}{\operatorname{Tr}}
\newcommand{\ii}{\mathrm{i}}
\newcommand{\dd}{\mathrm{d}}
\newcommand{\E}{\mathbb{E}}
\newcommand{\cE}{\mathcal{E}}
\newcommand{\cK}{\mathcal{K}}
\newcommand{\cT}{\mathcal{T}}
\newcommand{\cL}{\mathcal{L}}
\newcommand{\cJ}{\mathcal{J}}
\newcommand{\cO}{\mathcal{O}}
\newcommand{\bxi}{\bm{\xi}}
\newcommand{\br}{\bm r}
\newcommand{\bk}{\bm k}
\newcommand{\avg}[1]{\overline{#1}}
\newcommand{\cond}{\mathrm{cond}}
\newcommand{\Born}{\mathrm{B}}
\newcommand{\fixed}{\mathrm{fix}}
\newcommand{\ssub}{\mathrm{ss}}
\newcommand{\code}[1]{\texttt{#1}}

\begin{document}

\title{\texorpdfstring{Open-System / Monitored-System Response Theory\\
for a Class-A Adaptive Fermion Circuit}%
{Open-System / Monitored-System Response Theory for a Class-A Adaptive Fermion Circuit}}
\author{Internal theory and methods note}
\affiliation{Class-A fermionic adaptive-circuit project}
\date{August 18, 2026}

\begin{abstract}
This note consolidates the response-theory discussion surrounding chirality at
the monitored class-A domain wall.  Its first conclusion is economical: a local
number-phase impulse followed by a wall-resolved density measurement is already
sufficient to test directed propagation.  Fourier transformation along the
wall and in monitored circuit time should reveal a low-energy ridge near
$\omega=v_w k_y$, with the sign of $v_w$ reversing between the two oppositely
oriented walls.  This diagnoses chirality but is not, by itself, a quantized
Hall coefficient.  Its second conclusion is constructive and numerical.  For
perfect correction the Born-mean covariance obeys an exact affine replacement
map, so a time-dependent-flux pump can first be propagated deterministically,
with no outcome Monte Carlo, score estimator, or Keldysh calculation.  Physical
Born trajectories are then sampled only to measure record-to-record
fluctuations.  Event transfers and regional controller weights define an
inexpensive source-subtracted monitored-channel pump.  Because the reduced
circuit exchanges charge with feedback ancillas, a conventional electrical
Hall conductance requires one additional microscopic specification: compile
each local controller measurement into charge-conserving Givens gates, a local
system--ancilla reset, and the inverse gates, then count charge crossing a cut
gate by gate.  This produces a direct no-postselection numerical experiment in
the genuinely space--time-disordered circuit.  Generating-functional language
is retained only as optional analytic context.  Finally, we record the current
implementation status: the exact tangent formulas are already derived in the
manuscript appendix, whereas the present H2 production runner still uses
batched central $\pm\epsilon$ replays and should regard an analytic tangent as
a future upgrade.
\end{abstract}

\maketitle
\setcounter{tocdepth}{2}
\tableofcontents

\section{Purpose, scope, and bottom line}
\label{sec:purpose}

The original question was how a record-wise pattern can prove chirality, and
what a density-response theory means in a monitored trajectory ensemble.  The
discussion subsequently touched five related but inequivalent constructions:
individual trajectory response, trajectory-averaged response, Born-ensemble
response, density spectroscopy, and flux insertion.  Parts of the answer were
already distributed through the current manuscript and campaign documents, but
the complete logical chain was not available in one place.  This note is that
source of truth.

The minimum defensible program is:
\begin{enumerate}
  \item draw a physical stationary record $\bxi$ from the unperturbed Born
  process;
  \item apply an infinitesimal, charge-preserving local number-phase impulse at
  a source on one wall;
  \item propagate the exact covariance tangent along the same ordered record;
  \item retain both the conditional density derivative and the chronological
  derivative of the log Born probability;
  \item construct the fixed-record trajectory average and the full Born
  response separately;
  \item window the two walls separately and Fourier transform in the coordinate
  parallel to the wall and in physical circuit time;
  \item require an oriented low-energy ridge, opposite ridge slopes on opposite
  walls, wall localization, and a bulk or matched-trivial null result.
\end{enumerate}

No current operator is required for this claim.  A current-current response,
counting field, or adiabatic pump becomes necessary only for a stronger
transport or quantization statement.  The density susceptibility is a direct
chirality diagnostic because a one-dimensional chiral density mode has a
single signed dispersion at low energy.

\subsection{Where the material already lives}

Table~\ref{tab:repository-map} separates the compact manuscript statements,
the campaign authority, the current executable implementation, and the present
consolidation.  The older file
\path{00_WORKSPACE/CURRENT/prxq_draft/numerical_campaign.tex} contains
superseded H2/H3 sample counts and is not used as the operational authority.

\begin{table}[b]
\caption{Repository map.  Paths are relative to the current-workspace root.
``Formal'' means an exact derivation;
``executable'' means that the current production source actually evaluates the
stated estimator.}
\label{tab:repository-map}
\small
\begin{tabularx}{\textwidth}{@{}L{0.23\textwidth}Y L{0.16\textwidth}@{}}
\toprule
Location & Role & Status \\
\midrule
\path{prxq_draft/main.tex}, Sec.~``Chiral dynamics''
& Compact paper narrative for modular drift, full Born density response, and
fixed-record twist flow. & Formal, condensed \\
\path{prxq_draft/appendices.tex}
& Microscopic instrument, exact normalized tangent, chronological score,
Born-response identity, fixed-record holonomy, and ensemble distinctions. &
Formal \\
\path{experiment_review/numerical_campaign_legacy_working.tex}
& Audited H1--H3 run matrix, controls, acceptance gates, and legacy-evidence
classification. & Campaign authority \\
\path{final_production_ready_figure_scripts/_shared_src/fused_chirality_observables.py}
& H2 live-parent implementation; central differences the density and log
likelihood under common frozen records. & Executable, finite difference \\
\path{final_production_ready_figure_scripts/prior_designs/03_chirality_replay/}
& H3 representative frozen-record twist circle. & Executable \\
\path{Paper Campaign Summaries/mean_channel_lindblad_cpu_v2_hybrid_response/}
& Deterministic Gaussian mean-channel/Lindblad response with no Born records. &
Completed special case \\
This directory & Detailed synthesis, literature comparison, and exact-tangent
upgrade specification. & Canonical methods note \\
\bottomrule
\end{tabularx}
\end{table}

\section{The objects that must not be conflated}
\label{sec:taxonomy}

\subsection{Quantum instrument, record, and conditioned state}

Let a complete ordered record be
\begin{equation}
  \bxi=(e_1,m_1;e_2,m_2;\ldots;e_D,m_D),
\end{equation}
where $e_t$ includes the selected controller orbital, target occupation,
schedule information, and any exogenous random choice, while $m_t$ is the
intrinsic Born outcome.  The unnormalized branch map is
\begin{equation}
  \cK_{\bxi,u}(\rho_0)
  =K_{D,u}\cdots K_{1,u}\rho_0
   K_{1,u}^{\dagger}\cdots K_{D,u}^{\dagger} .
\end{equation}
For perturbation parameter $u$, its trace and normalized state are
\begin{align}
  p_u(\bxi)&=\Tr\cK_{\bxi,u}(\rho_0),\label{eq:path-probability}\\
  \rho_{u,\bxi}&=\frac{\cK_{\bxi,u}(\rho_0)}{p_u(\bxi)} .
  \label{eq:conditional-state}
\end{align}
Equation~\eqref{eq:path-probability} is the many-body Born probability of the
entire record.  A Born-sampled Monte Carlo ensemble already draws each $\bxi$
with this probability.  Its trajectories therefore enter ordinary sample
means with equal numerical weight; multiplying each sample by its stored
$p(\bxi)$ a second time would double-weight the ensemble.

Quantum-trajectory unravelings and their equivalence to master-equation
averages for linear observables are standard
\cite{DalibardCastinMolmer1992,CarischRomitoZilberberg2023}.  What matters here
is that the adaptive record also changes the future normalized state and hence
all later conditional probabilities.

\subsection{Three meanings of response}

Let
\begin{equation}
  O_{u,\bxi}=\Tr[O\rho_{u,\bxi}]
\end{equation}
be a trajectory observable.  There are three useful response objects.

First, the \emph{trajectory-conditioned response} holds the complete record
fixed:
\begin{equation}
  \chi_{\bxi}^{\cond}
  =\left.\partial_u O_{u,\bxi}\right|_{u=0}.
  \label{eq:conditional-response}
\end{equation}
This is a counterfactual branch derivative.  It asks what the normalized state
would do under a weak perturbation if precisely the same sequence of outcomes
were imposed.

Second, the \emph{fixed-record trajectory average} is
\begin{equation}
  \chi^{\fixed}
  =\E_{p_0}\!\left[\chi_{\bxi}^{\cond}\right].
  \label{eq:fixed-average}
\end{equation}
It samples ordinary unperturbed records but does not differentiate their
probabilities.  We use ``fixed-record'' or ``quenched conditional average'' to
remove the ambiguity in the informal phrase ``trajectory-averaged response.''

Third, the \emph{Born-ensemble response} differentiates the operational mean
of the perturbed experiment:
\begin{equation}
  \chi^{\Born}
  =\left.\partial_u\sum_{\bxi}p_u(\bxi)O_{u,\bxi}\right|_{u=0}.
  \label{eq:born-response-definition}
\end{equation}
Defining the path score
\begin{equation}
  S_{\bxi}=\left.\partial_u\log p_u(\bxi)\right|_{u=0},
  \label{eq:path-score}
\end{equation}
ordinary differentiation gives the exact identity
\begin{equation}
  \boxed{
  \chi^{\Born}
  =\E_{p_0}\!\left[
       \chi_{\bxi}^{\cond}+O_{0,\bxi}S_{\bxi}
    \right].}
  \label{eq:score-identity}
\end{equation}
Normalization implies $\E_{p_0}[S_{\bxi}]=0$, so a variance-reduced form is
\begin{equation}
  \chi^{\Born}
  =\chi^{\fixed}
   +\operatorname{Cov}_{p_0}(O_{0,\bxi},S_{\bxi}).
  \label{eq:score-covariance}
\end{equation}
The score is familiar from inference on continuously monitored systems
\cite{GammelmarkMolmer2013}, and recent work derives response bounds for quantum
trajectory observables \cite{Vu2025}.  Equations~\eqref{eq:score-identity} and
\eqref{eq:score-covariance} are nevertheless elementary likelihood-ratio
identities: their role here is to identify the physical estimator unambiguously.

If separate experiments are run at $u=+\epsilon$ and $u=-\epsilon$, with
records resampled from $p_{+\epsilon}$ and $p_{-\epsilon}$, the derivative of
those two sample means estimates $\chi^{\Born}$.  By contrast, replaying the
same baseline records at both signs estimates the two terms of
Eq.~\eqref{eq:score-identity} and exposes their decomposition.

\subsection{Born ensemble versus unconditional state}

The unconditional state is
\begin{equation}
  \bar\rho_u=\sum_{\bxi}p_u(\bxi)\rho_{u,\bxi}.
  \label{eq:unconditional-state}
\end{equation}
For any linear quantum observable,
\begin{equation}
  \sum_{\bxi}p_u(\bxi)O_{u,\bxi}
  =\Tr(O\bar\rho_u).
  \label{eq:linear-equivalence}
\end{equation}
Consequently the full Born density response can, in principle, be calculated
from the differentiated unconditional instrument.  Equation~\eqref{eq:score-identity}
is an unraveling of that same derivative.  It is valuable because it separates
what happens to a typical conditional state from what happens because the
record distribution is reweighted.

This equivalence fails for nonlinear trajectory observables:
\begin{equation}
  \E[S(\rho_{\bxi})]\ne S(\bar\rho),\qquad
  \E[|C_{ij,\bxi}|^2]\ne|\bar C_{ij}|^2.
\end{equation}
Entanglement, trajectory topology, purity, and a per-record fitted velocity must
therefore be formed trajectory by trajectory before averaging.  An adaptive
schedule with finite classical memory may also require an enlarged
classical--quantum channel rather than a system-only one; recent feedback full
counting statistics makes this construction explicit
\cite{RosalFiusaPottsLandi2026}.

\begin{table}[t]
\caption{Response taxonomy used throughout this note.}
\label{tab:response-taxonomy}
\small
\begin{tabularx}{\textwidth}{@{}L{0.20\textwidth}L{0.28\textwidth}Y@{}}
\toprule
Object & Formula & Operational statement \\
\midrule
Conditional response & $\chi_{\bxi}^{\cond}=\partial_uO_{u,\bxi}|_0$ &
Same complete record is imposed before and after the perturbation. \\
Fixed-record average & $\chi^{\fixed}=\E_{p_0}[\chi_{\bxi}^{\cond}]$ &
Average conditional sensitivity of ordinary unperturbed trajectories. \\
Score contribution & $\E_{p_0}[(O-\bar O)S]$ &
Correction from the perturbation-induced redistribution of records. \\
Born response & $\partial_u\E_{p_u}[O_u]|_0$ &
Repeated perturbed experiment with outcomes sampled from their perturbed Born
probabilities. \\
Unconditional response & $\partial_u\Tr(O\bar\rho_u)|_0$ &
Identical to the Born response for a linear $O$, provided the complete
instrument and classical feedback state are included. \\
\bottomrule
\end{tabularx}
\end{table}

\section{Exact Gaussian trajectory response}
\label{sec:exact-tangent}

\subsection{Covariance and sign convention}

For the canonical engine we use
\begin{equation}
  C_{ij}=\langle c_i^\dagger c_j\rangle,
  \qquad G=2C-\Id,
  \label{eq:covariance-convention}
\end{equation}
so $G$ is Hermitian with spectrum in $[-1,1]$.  There is a persistent transpose
convention hazard in the wider repository: the alternative one-particle density
matrix $\widetilde C_{ij}=\langle c_j^\dagger c_i\rangle=C_{ji}$ transforms by
the adjoint single-particle convention.  All derivative signs must therefore be
checked against the implemented matrix transformation, rather than inferred
from a shorthand phrase such as $C=\langle c^\dagger c\rangle$.

The current H2 code implements the source as
\begin{equation}
  G_\epsilon=D_\epsilon G D_\epsilon^\dagger,
  \qquad D_\epsilon=e^{-\ii\epsilon P_0},
  \label{eq:implemented-kick}
\end{equation}
and hence
\begin{equation}
  X_0\equiv\left.\partial_\epsilon G_\epsilon\right|_0
  =-\ii[P_0,G].
  \label{eq:initial-tangent}
\end{equation}
Under Eq.~\eqref{eq:covariance-convention}, the many-body unitary producing
Eq.~\eqref{eq:implemented-kick} is $U_\epsilon=e^{+\ii\epsilon n_0}$.
Writing instead the conventional pulse $e^{-\ii\epsilon n_0}$ replaces
$\epsilon\mapsto-\epsilon$ and flips the overall response sign.  It does not
change the signed slope of the $(k,\omega)$ ridge or the reversal between wall
orientations.  Equation~\eqref{eq:implemented-kick} is the authoritative
convention for comparisons with the executable H2 product.

The impulse preserves total particle number and initially changes only
coherences.  Since all local number projectors commute, the instantaneous
density response vanishes at the source time; the subsequent monitored
dynamics converts the coherence tangent into a signed density pattern.  This
is fundamentally different from resetting a cell from empty to occupied,
which injects a finite charge and is not a linear-response estimator.

\subsection{One normalized measurement-and-reset event}

Let $P=|\chi\rangle\langle\chi|$ be a normalized measured orbital,
$Q=\Id-P$, $m\in\{0,1\}$ its outcome, and
$\eta_m=2m-1$.  Let $s\in\{0,1\}$ denote the controller's target occupation.
The Born probability is
\begin{equation}
  p_m(G)=\frac{1+\eta_m g}{2},
  \qquad g=\Tr(PG)=\langle\chi|G|\chi\rangle.
  \label{eq:binary-born}
\end{equation}
For a nonzero-probability branch, projective measurement followed by perfect
reset acts as
\begin{equation}
  \boxed{
  \Phi_{m|s}(G)
  =Q(\Id+\eta_m GP)^{-1}GQ+\eta_sP.}
  \label{eq:conditional-map}
\end{equation}
The apparent inverse is only a rank-one resolvent.  Define
\begin{equation}
  L_m(G)=Q(\Id+\eta_m GP)^{-1}
  =Q-\frac{\eta_m QGP}{1+\eta_m g}.
  \label{eq:left-jacobian}
\end{equation}
At fixed $P,m,s$, exact differentiation gives the congruence
\begin{equation}
  \boxed{
  D_G\Phi_{m|s}\big|_G[X]
  =L_m(G)X L_m(G)^\dagger.}
  \label{eq:event-tangent}
\end{equation}
The target $s$ does not enter the derivative because the reset covariance is a
constant.  Exact projectors make $L_m$ singular: $L_mP=0$.  This rank loss is a
physical part of the normalized response and must not be replaced by a
pseudoinverse.

The current manuscript appendix uses the equivalent right-factor convention
$J_m=L_m^\dagger$ and writes $D\Phi[X]=J_m^\dagger XJ_m$.  We use the left
factor because chronological propagation then has its familiar order.

\subsection{Chronological tangent and general parameter dependence}

Along a fixed record, let $L_e$ be Eq.~\eqref{eq:left-jacobian} evaluated on
the state immediately before event $e$.  If the perturbation is only the
initial impulse and does not change later projectors, then
\begin{align}
  X_{e+1}&=L_eX_eL_e^\dagger,\label{eq:tangent-recursion}\\
  X_D&=A_{D:0}X_0A_{D:0}^\dagger,
  \qquad A_{D:0}=L_{D-1}\cdots L_1L_0.
  \label{eq:tangent-product}
\end{align}
This is an exact derivative of the normalized conditioned trajectory, not a
derivative of an unnormalized Kraus product and not a finite difference.

For a parameter that also changes the measured orbital, write
$P=P(u)$, $Q=\Id-P$, $A=\Id+\eta_m GP$, and $X=\partial_uG$.  Direct
Fr\'echet differentiation of Eq.~\eqref{eq:conditional-map} gives
\begin{multline}
  \partial_u\Phi_{m|s}
  =\dot Q A^{-1}GQ
   +Q\left[-A^{-1}\dot A A^{-1}G+A^{-1}X\right]Q\\
   +QA^{-1}G\dot Q+\eta_s\dot P,
  \qquad
  \dot A=\eta_m(XP+G\dot P).
  \label{eq:general-frechet}
\end{multline}
Equation~\eqref{eq:general-frechet} is the appropriate starting point for a
gauge-field or twist tangent, where the controller orbitals themselves carry
Peierls phases.  H2 needs only the simpler homogeneous recursion
Eq.~\eqref{eq:tangent-recursion}.

\subsection{Exact chronological score}

The record probability factorizes conditionally:
\begin{equation}
  p_u(\bxi)=\prod_e p_{m_e}\!\left(G_e(u),P_e(u)\right).
  \label{eq:record-product}
\end{equation}
The score is therefore accumulated event by event,
\begin{align}
  S_{e+1}&=S_e+s_e,\qquad S_0=0,\label{eq:score-recursion}\\
  s_e&=\frac{\eta_{m_e}\dot g_e}{1+\eta_{m_e}g_e},
  \label{eq:binary-score}\\
  \dot g_e&=\Tr(P_eX_e)+\Tr(\dot P_eG_e).
  \label{eq:g-derivative}
\end{align}
For the H2 initial density impulse, $\dot P_e=0$ after the kick.  The tangent
$X_e$ contains the complete effect of all earlier nonlinear normalized
updates.  The score cannot be obtained by differentiating every probability
at the initial covariance.

For a probe cell $r$ with two-orbital projector $P_r$, the density and its
conditional derivative are
\begin{align}
  n_{r,\bxi}(t)&=\Tr(P_rC_{\bxi}(t)),\label{eq:cell-density}\\
  \chi_{r,\bxi}^{\cond}(t;t_0)
  &=\frac12\Tr[P_rX_{\bxi}(t;t_0)].
  \label{eq:density-tangent}
\end{align}
The full wall-density response is then
\begin{multline}
  \chi_{r}^{\Born}(t;t_0)
  =\E_{p_0}\!\left[\frac12\Tr(P_rX_{\bxi}(t;t_0))\right]\\
  +\E_{p_0}\!\left[
  \bigl(n_{r,\bxi}(t)-\bar n_r(t)\bigr)S_{\bxi}(t;t_0)
  \right].
  \label{eq:full-density-response}
\end{multline}
Equation~\eqref{eq:full-density-response} is the central H2 formula.  Its first
line is the fixed-record trajectory average; its second is the Born-score
correction.

\subsection{Finite differences as validation, not definition}

The current production runner forms
\begin{align}
  \partial_\epsilon n_{\bxi}
  &\simeq\frac{n_{+\epsilon,\bxi}-n_{-\epsilon,\bxi}}{2\epsilon},\nonumber\\
  S_{\bxi}
  &\simeq\frac{\log p_{+\epsilon}(\bxi)-
                  \log p_{-\epsilon}(\bxi)}{2\epsilon},
  \label{eq:current-finite-difference}
\end{align}
using common frozen schedules and outcomes.  It uses $\epsilon=10^{-3}$ and
retains $5\times10^{-4}$ and $2\times10^{-3}$ as validation scales.  This is a
legitimate central-difference estimator and its GPU batching is efficient, but
it is not the analytic tangent of Eqs.~\eqref{eq:tangent-recursion} and
\eqref{eq:score-recursion}.  The exact tangent removes step-size bias,
cancellation error, and the need for two full replays per source.  The central
difference should remain as a small-system parity test over an $\epsilon$
convergence window.

\section{Density response as a chirality diagnostic}
\label{sec:density-chirality}

\subsection{Response, correlation, and fluctuation--dissipation}

Three density objects are often called a ``density correlator'' in informal
discussion:
\begin{align}
  C_{nn}(r,t)&=\langle n(r,t)n(0,0)\rangle_c,
  &&\text{spontaneous correlation},\nonumber\\
  \chi^R_{nn}(r,t)&=-\ii\theta(t)
  \langle[n(r,t),n(0,0)]\rangle,
  &&\text{closed equilibrium response},\nonumber\\
  \chi_{nn}(r,t)&=\left.\frac{\delta\langle n(r,t)\rangle_V}
  {\delta V(0,0)}\right|_{V=0},
  &&\text{operational susceptibility}.
  \label{eq:three-density-objects}
\end{align}
In a closed equilibrium system, Kubo theory relates the last two objects
\cite{Kubo1957}.  A monitored nonequilibrium trajectory generally does not obey
the equilibrium fluctuation--dissipation theorem.  The dissipator, normalized
backaction, and record probabilities all respond.  Open-system NESS response
theory explicitly warns that the answer need not reduce to a simple
system-only two-point correlator
\cite{KonopikLutz2019,LevyRabaniLimmer2021}.  For this project the safest
definition is therefore the operational derivative in the last line of
Eq.~\eqref{eq:three-density-objects}, evaluated exactly by
Eq.~\eqref{eq:full-density-response}.

Static density correlations can demonstrate a critical wall and determine a
scaling exponent, but they do not determine the propagation direction.  The
retarded density response contains that missing arrow of motion.

\subsection{Real-space pattern}

Let $x_w$ denote one domain wall and let $f_w(x)$ be a narrow, preregistered
transverse window.  For a source at $(x_w,y_0)$ define
\begin{equation}
  \chi_w(r,t)
  =\sum_x f_w(x)\,
  \chi^{\Born}_{(x,y_0+r),(x_w,y_0)}(t;t_0),
  \qquad r=y-y_0\pmod{N_y}.
  \label{eq:wall-window-response}
\end{equation}
A chiral response should satisfy the following conjunction:
\begin{enumerate}
  \item the response remains concentrated in $x$ near the chosen wall;
  \item its signed or squared weight moves predominantly in one $y$ direction;
  \item the propagation direction reverses on the oppositely oriented wall;
  \item a source in the bulk lacks the isolated directional wall feature;
  \item a matched trivial or wall-off circuit lacks the same low-energy ridge;
  \item the sign is stable against the source, late time origin, wall-window
  width, and fit window.
\end{enumerate}
The density-phase impulse has zero net injected charge, so $\sum_r\chi_w(r,t)$
can be small or vanish and the response can be a signed dipole.  A centroid
must therefore use a declared nonnegative weight such as $|\chi|$ or
$|\chi|^2$, while the signed response itself is retained for physics.  A useful
real-space directional statistic is
\begin{equation}
  \mathcal C_w(t)=
  \frac{W_{w,+}(t)-W_{w,-}(t)}{W_{w,+}(t)+W_{w,-}(t)},
  \label{eq:real-space-chirality}
\end{equation}
where $W_{w,\pm}$ are response weights on the two directed arcs before
wraparound.

\subsection{Why a Fourier ridge is chiral}

For a single low-energy chiral density branch, chiral Luttinger theory gives,
up to normalization and Fourier conventions,
\begin{equation}
  \chi^R_{nn}(k,\omega)
  \propto\frac{k}{\omega-vk+\ii0^+}.
  \label{eq:chiral-pole}
\end{equation}
This is the equilibrium reference prediction
\cite{Wen1990}.  Density-coupled spectroscopy and real-time charge propagation
have been used to expose chiral edge dispersions in closed lattice models
\cite{GoldmanBeugnonGerbier2012,DongGrushinMotrukPollmann2018}.

Monitoring replaces an infinitely sharp pole by a generally broadened mode,
\begin{equation}
  \chi_w(k,\omega)
  \sim\frac{Z_w(k)}{\omega-v_wk+\ii\Gamma_w(k)}
       +\chi_w^{\mathrm{reg}}(k,\omega),
  \label{eq:broadened-pole}
\end{equation}
or, for a stroboscopic channel eigenmode,
\begin{equation}
  \lambda_w(k)=e^{-\gamma_w(k)+\ii\omega_w(k)}.
  \label{eq:channel-dispersion}
\end{equation}
The peak location measures $\omega_w(k)$, while its linewidth reflects
$\gamma_w(k)$.  At low momentum one expects
\begin{equation}
  \omega_w(k)\simeq v_w k,
  \qquad v_{w_1}v_{w_2}<0.
  \label{eq:opposite-velocities}
\end{equation}
The density operator probes particle--hole density excitations, not the bare
single-particle edge band.  For a nearly linear edge they share the same
characteristic velocity; band curvature, monitoring, and bulk leakage broaden
the response into a continuum.

The lattice edge spectrum remains periodic over the Brillouin zone.  Its
topological feature is spectral flow connecting bulk continua, not literal
nonperiodicity of an edge energy function.  Equation~\eqref{eq:chiral-pole} is
a low-energy description near the crossing.

\subsection{Discrete Fourier protocol and symmetry}

After translating every source to $r=0$ within its parent trajectory, use
\begin{equation}
  \chi_w(k_j,\omega_\ell)
  =\sum_{t=0}^{T_R}\sum_{r=0}^{N_y-1}
  w_t\,e^{\ii\omega_\ell t-\ii k_jr}\chi_w(r,t),
  \label{eq:discrete-fourier}
\end{equation}
with $k_j=2\pi j/N_y$ and a declared causal time taper $w_t$.  Ridge fitting
uses the native frequency resolution; zero padding is a visualization choice.
Because circuit time is discrete, quasifrequency is periodic modulo $2\pi$ per
cycle, so the fit must remain inside a low-frequency window away from aliasing
and short-time lattice structure.

For a real time-domain response,
\begin{equation}
  \chi_w(-k,-\omega)=\chi_w(k,\omega)^*.
  \label{eq:reality-symmetry}
\end{equation}
The unavoidable $(-k,-\omega)$ image is not a counterpropagating partner.  One
may restrict to $\omega>0$ and compare the two momentum signs, or use a wedge
asymmetry
\begin{equation}
  \mathcal A_w=
  \frac{\sum_{(k,\omega)\in\mathcal W_+}W_w(k,\omega)
       -\sum_{(k,\omega)\in\mathcal W_-}W_w(k,\omega)}
       {\sum_{(k,\omega)\in\mathcal W_+}W_w(k,\omega)
       +\sum_{(k,\omega)\in\mathcal W_-}W_w(k,\omega)},
  \label{eq:wedge-asymmetry}
\end{equation}
where $\mathcal W_\pm$ are matched low-energy regions with
$\operatorname{sgn}(\omega k)=\pm1$.  $W_w$ may be $|\chi_w|^2$ for robust
ridge localization, but it must not be called the retarded spectral function.
When phase information and conventions are controlled,
$-\operatorname{Im}\chi^R$ carries the sharper signed spectral meaning.

The two physical walls must never be summed before this analysis: their
opposite velocities can cancel.  Source positions and time origins are
correlated variance-reduction samples inside a trajectory.  They do not
increase the number of statistically independent trajectories.

\section{What record-wise evidence can and cannot prove}
\label{sec:recordwise}

\subsection{Distribution of conditioned responses}

For every parent trajectory, H2 can report a record-conditioned real-space
map, Fourier ridge velocity $v_{w,\bxi}$, and asymmetry
$\mathcal A_{w,\bxi}$.  The scientifically meaningful questions are then:
\begin{itemize}
  \item Does the distribution concentrate at a nonzero sign?
  \item Does the sign reverse between the paired wall orientations?
  \item How large is the score correction relative to the fixed-record mean?
  \item Does the wrong-chirality tail shrink with $N_y$?
  \item Are outliers associated with small Born denominators, leakage, or
  nonstationarity?
\end{itemize}

Finite Monte Carlo data cannot prove that \emph{every} possible record is
chiral.  A complete record space is exponentially large and can include rare
branches with anomalous response.  The defensible trajectory-level statement
is an almost-sure or high-probability claim supported by the distribution,
quantiles, sign fraction, finite-size evolution, and preregistered failure
counts.  The Born mean can be chiral even if some conditioned records are not;
conversely, a chiral fixed-record mean can be altered or canceled by the score
term.

\subsection{H3 frozen-record flux flow is a different record-wise question}

H3 freezes one complete physical record sampled at zero twist and replays that
same record for $0\le\phi\le2\pi$.  Conceptually one seeks a wall-resolved
entanglement flow.  The current executable restricts $C=(G+\Id)/2$ to one
half-torus containing all $x$ and $0\le y<N_y/2$, selects the eigenmodes nearest
occupation $1/2$, and assigns them to nominal wall windows by their $x$ weight.
Let $\nu_a(\phi)$ be the resulting overlap-tracked covariance eigenvalues and
define the entanglement levels
\begin{equation}
  \varepsilon_a(\phi)=\log\frac{1-\nu_a(\phi)}{\nu_a(\phi)}.
  \label{eq:entanglement-level}
\end{equation}
Oriented crossings of $\varepsilon=0$ diagnose branchwise spectral flow.  The
equilibrium analogy is Laughlin flux insertion on a cylinder
\cite{Laughlin1981}: a $2\pi$ twist shifts allowed $k_y$ values by one momentum
step.  Bulk bands return to themselves as sets, while a chiral edge branch
undergoes a level permutation through the gap.  Oppositely oriented interfaces
carry opposite permutations, so the total flow of the periodic two-wall
sample vanishes.

The record must remain fixed because Berry holonomy or spectral flow requires
a continuous family.  Independently resampling $\bxi(\phi)$ at every twist
produces unrelated states, and no amount of eigenvalue sorting creates a
physical eigenbundle.  The family is defined only where its branch probability
$p_{\bxi}(\phi)$ remains nonzero and its relevant rank is constant.

The approved production H3 scope is one preregistered interface record and one
matched-trivial record at $20\times24$, each on a 65-point closed circle.  The
seam gauge places the nontrivial link between $y=23$ and $y=0$.  H3 is therefore
a representative branch-level diagnostic.  It cannot establish
the fraction of Born records with spectral flow, typicality, or ensemble
quantization.  H3 complements H2 but does not replace it.

\section{Three clocks for chirality}
\label{sec:three-clocks}

The project now has three primary chirality routes.  Their agreement would be
powerful precisely because their clocks and claims differ.

\begin{table}[t]
\caption{Three chirality clocks and their claim boundaries.}
\label{tab:three-clocks}
\small
\begin{tabularx}{\textwidth}{@{}L{0.12\textwidth}L{0.18\textwidth}L{0.27\textwidth}Y@{}}
\toprule
Route & Clock & Positive signature & Does not establish \\
\midrule
H1 modular drift & Modular time generated by $K_A=-\log\rho_A$ & Opposite
packet drifts on opposite walls, formed trajectory by trajectory & Physical
circuit-time transport or a Hall coefficient \\
H2 density response & Monitored circuit time & Wall-localized one-way
$\chi(k,\omega)$ ridge; conditional and score pieces; orientation reversal &
Quantized transport by itself \\
H3 frozen twist & Auxiliary twist angle $\phi$ & Continuous wall-resolved
entanglement spectral flow for the same record & Born-ensemble frequency,
typicality, or a measured DC current \\
\bottomrule
\end{tabularx}
\end{table}

H1 asks whether the prepared state has a handed modular generator.  H2 asks
whether the physical monitored dynamics carries a causal density disturbance
in one direction.  H3 asks whether a continuous conditional branch has
oriented flux spectral flow.  Static entanglement scaling fixes critical
content such as $c_L+c_R$, but it does not determine the propagation
orientation $c_R-c_L$; at least one orientation-sensitive route is necessary.

\section{Steady-state and open-system response}
\label{sec:open-response}

\subsection{A steady state is not a response theory}

It is tempting to ``take the steady state and look at it.''  A stationary
covariance can reveal an occupation or purity gap and a topological marker, but
it does not determine how the system responds.  The same $\rho_{\ssub}$ can be
stationary under distinct generators with different relaxation poles,
currents, and susceptibilities.  One needs the steady state \emph{and} the
dynamical generator or quantum instrument, including how the perturbation
couples to it.  This distinction is central in classifications of quadratic
Lindbladians, where steady-state topology and damping-spectrum topology are
not interchangeable \cite{LieuMcGinleyCooper2020}.

For a differentiable discrete unconditional channel
\begin{equation}
  \rho_{n+1}=\cE_u(\rho_n),
  \qquad \cE_0(\rho_{\ssub})=\rho_{\ssub},
\end{equation}
the tangent obeys
\begin{equation}
  X_{n+1}=\cE_0(X_n)+(\partial_u\cE_u)_0(\rho_{\ssub})u_n.
  \label{eq:channel-linearization}
\end{equation}
For harmonic $u_n=u_\omega z^n$, the stationary response contains the channel
resolvent
\begin{equation}
  X(\omega)=
  \left(z\Id-\cE_0\right)^{-1}
  (\partial_u\cE_u)_0(\rho_{\ssub})u_\omega,
  \qquad z=e^{-\ii\omega}.
  \label{eq:discrete-resolvent}
\end{equation}
The eigenvalues of $\cE_0$ determine resonance positions and damping.

For a continuous Markovian generator,
\begin{equation}
  \dot\rho=\cL_u\rho,
\end{equation}
one analogously obtains
\begin{equation}
  \dot X=\cL_0X+(\partial_u\cL_u)_0\rho_{\ssub}\,u(t),
  \qquad
  X(\omega)=(\ii\omega-\cL_0)^{-1}
  (\partial_u\cL_u)_0\rho_{\ssub}\,u(\omega),
  \label{eq:lindblad-resolvent}
\end{equation}
for the stated Fourier convention.  Lindblad and NESS response theories are
well established
\cite{AvronFraasGraf2012,VenutiZanardi2016,AlbertBradlynFraasJiang2016,KonopikLutz2019}.
Gaussian-channel response can also close at the covariance level
\cite{MehboudiParrondoAcin2019}.  These are unconditional theories; the exact
trajectory tangent plus score is their record-resolved realization for the
present instrument.

\subsection{Deterministic mean-channel special case}

The completed CPU mean-channel campaign studies
\begin{equation}
  \dot C=-DC-CD+S,
  \qquad D C_{\ssub}+C_{\ssub}D=S.
  \label{eq:mean-channel}
\end{equation}
For a density-phase impulse with
$\delta C(0)=-\ii[P_0,C_{\ssub}]$, its exact response is
\begin{equation}
  \delta C(t)=-\ii e^{-Dt}[P_0,C_{\ssub}]e^{-Dt}.
  \label{eq:mean-channel-response}
\end{equation}
This supplies a useful analytic calibration and a controlled special case of
Eqs.~\eqref{eq:lindblad-resolvent} and \eqref{eq:broadened-pole}.  It contains
no sampled measurement records, no nonlinear normalized branch, and no score
term.  It must not be presented as an H2 Born-trajectory result.  Conversely,
a gapped deterministic mean channel does not rule out slow or chiral
conditioned trajectory structure.

\section{No-postselection Hall transport in the space--time-disordered circuit}
\label{sec:quantization}

\subsection{Density chirality is already a coherent target}

H2's immediate claim is directionality: a localized, charge-preserving
perturbation produces a wall-confined response whose dispersion has only one
low-energy velocity sign, reversed on the other wall.  This is sufficient to
show a chiral domain-wall mode.  It does not assign a quantized transport
coefficient.

Open-system topology and transport are not generically locked.  A topological
steady covariance can coexist with nonquantized Hall response, and special
conditions are required for quantization
\cite{ShavitGoldstein2020,AlbertBradlynFraasJiang2016}.  Finite-range
dissipative state preparation also obeys nontrivial locality constraints
\cite{Goldstein2019}.  The manuscript should therefore keep ``state
topology,'' ``chiral density response,'' and ``quantized transport'' as three
separate claims.

\subsection{Primary microscopic numerical definition}

No field theory is needed to define or estimate the response.  The basic data
are only the covariance trajectory, the ordered measurement record, and the
controller orbitals already used by the simulator.

Partition the periodic $x$ direction into two complementary regions $R_L$ and
$R_R$, each containing one interface, with the two artificial partition cuts
placed deep in the topological and trivial bulks.  For the production
$N_x=20$ explicit-interface geometry, a convenient preregistered choice is
\begin{equation}
  R_L:\ x=0,\ldots,9,
  \qquad
  R_R:\ x=10,\ldots,19.
  \label{eq:hall-regions}
\end{equation}
The physical interfaces lie on bonds $(4,5)$ and $(15,16)$, whereas the
partition cuts lie near the midpoints of the topological and trivial sectors.
For any one-particle region projector $R_a$, $a=L,R$, define the trajectory
particle number
\begin{equation}
  N_{a,\bxi}(c)
  =\frac12\Tr\!\left[R_a(\Id+G_{\bxi}(c))\right].
  \label{eq:regional-particle-number}
\end{equation}

At event $e$, the normalized controller orbital at the instantaneous twist is
$|\chi_e(\phi_e)\rangle$, with
$P_e(\phi_e)=|\chi_e(\phi_e)\rangle\langle\chi_e(\phi_e)|$.  If the Born
outcome is $m_e\in\{0,1\}$ and the perfect-correction target is
$s_e\in\{0,1\}$, the feedback ancilla changes the physical covariance by
$(s_e-m_e)P_e$.  Its contribution to region $a$ is therefore the exactly
known number
\begin{equation}
  A_{a,\bxi}(c)
  =\sum_{e\le c}(s_e-m_e)
    \Tr\!\left[R_aP_e(\phi_e)\right].
  \label{eq:microscopic-ancilla-source}
\end{equation}
The source-subtracted regional balance is
\begin{equation}
  q_{a,\bxi}(c)
  =N_{a,\bxi}(c)-N_{a,\bxi}(0)-A_{a,\bxi}(c),
  \label{eq:microscopic-region-flow}
\end{equation}
in units of electron number.  Finally define the oriented reduced transverse transfer
between the two walls by
\begin{equation}
  q_{x,\bxi}(c)
  =\frac12\left[q_{R,\bxi}(c)-q_{L,\bxi}(c)\right].
  \label{eq:microscopic-transverse-transfer}
\end{equation}
Because $R_L+R_R=\Id$, exact perfect-correction arithmetic requires
$q_L+q_R=0$ up to numerical roundoff.  Equation
\eqref{eq:microscopic-transverse-transfer} consequently uses both regions as
an internal continuity check rather than relying on an assumed current
operator.  Record by record, this reduced quantity also contains the
measurement-innovation change of a conditioned regional expectation.  It is
therefore a precise and inexpensive monitored-channel pump estimator, but not
yet a literal electrical current.  This innovation cancels in the ordinary
Born mean of a local compiled measurement; it matters critically for
record-level current distributions.

Everything in Eqs.~\eqref{eq:regional-particle-number}--
\eqref{eq:microscopic-transverse-transfer} is available microscopically.  The
cycle observer supplies $G(c)$; the record observer supplies $s_e-m_e$, the
site, and the channel; and the twisted controller spinors supply
\begin{equation}
  w_a(j,\nu;\phi)
  =\Tr\!\left[R_aP_{j,\nu}(\phi)\right]
  =\sum_{i\in R_a}|\chi_{j,\nu,i}(\phi)|^2.
  \label{eq:regional-controller-weight}
\end{equation}
Thus the observable can be streamed as two scalar arrays per trajectory and
cycle.  No covariance history, Keldysh functional, or guessed bond-current
operator is required.

\subsection{Exact deterministic Born mean}

For the canonical perfect-correction instrument, the Born mean of every linear
covariance observable closes exactly.  This removes even the need for
trajectory Monte Carlo when the first target is the ensemble-mean Hall
coefficient.

For one normalized rank-one mode $P_e$, let $Q_e=\Id-P_e$ and let its target
occupation be $s_e\in\{0,1\}$.  Averaging the two normalized measurement and
reset branches with their actual Born probabilities gives
\begin{equation}
  \overline C_{e+1}
  =Q_e(\phi_e)\overline C_eQ_e(\phi_e)
   +s_eP_e(\phi_e).
  \label{eq:exact-born-mean-replacement}
\end{equation}
The nonlinear branch denominators cancel in the weighted sum.  In the
$G=2C-\Id$ convention this is
\begin{equation}
  \overline G_{e+1}
  =Q_e\overline G_eQ_e+(2s_e-1)P_e.
  \label{eq:exact-born-mean-G}
\end{equation}
Equation~\eqref{eq:exact-born-mean-G} is already the public
\code{mean\_replacement=True} path of the canonical GPU engine.  For a fixed
exogenous site schedule it is the exact Born average over all intrinsic
measurement words, not an approximation.  Random schedules can be sampled in
parallel; each GPU batch member then represents one schedule realization with
its Born outcomes summed analytically.

The expected feedback source is equally microscopic.  Immediately before
event $e$,
\begin{equation}
  \overline y_e
  =\E[s_e-m_e]
  =s_e-\Tr(P_e\overline C_e),
  \label{eq:exact-mean-transfer}
\end{equation}
so that
\begin{equation}
  \overline A_a(c)
  =\sum_{e\le c}
  \left[s_e-\Tr(P_e\overline C_e)\right]
  \Tr(R_aP_e).
  \label{eq:exact-mean-regional-source}
\end{equation}
Propagating Eqs.~\eqref{eq:exact-born-mean-replacement} and
\eqref{eq:exact-mean-regional-source} therefore gives the exact Born-mean
version of Eqs.~\eqref{eq:microscopic-region-flow} and
\eqref{eq:microscopic-transverse-transfer}.  The only sampling error is from an
optional random schedule ensemble, not from Born outcomes.

An analytic microscopic tangent is also finite-dimensional.  With dots
denoting differentiation with respect to the applied flux,
\begin{equation}
  \dot{\overline C}_{e+1}
  =\dot Q_e\overline C_eQ_e
   +Q_e\dot{\overline C}_eQ_e
   +Q_e\overline C_e\dot Q_e
   +s_e\dot P_e,
  \qquad \dot Q_e=-\dot P_e.
  \label{eq:exact-mean-flux-tangent}
\end{equation}
This is a purely numerical matrix recursion through the microscopic event
sequence.  It computes the linear Born response without a finite difference,
without a likelihood score, and without a correlation function.  The score is
needed only when one insists on resolving the derivative trajectory by
trajectory before averaging; Eq.~\eqref{eq:exact-born-mean-replacement} has
already performed that average exactly.

The most efficient campaign hierarchy is consequently:
\begin{enumerate}
  \item propagate Eq.~\eqref{eq:exact-born-mean-replacement} under the flux
  protocol to establish whether the Born mean has a plateau;
  \item validate its flux derivative either by
  Eq.~\eqref{eq:exact-mean-flux-tangent} or paired finite differences;
  \item only if the mean result is promising, run physical Born trajectories
  to measure variance, quantiles, and record concentration.
\end{enumerate}
The current mean-replacement path needs one lightweight event observer to emit
Eq.~\eqref{eq:exact-mean-transfer}; this is smaller than implementing a new
response formalism.

\subsection{Microscopic electrical current from a local gate compilation}

A literal current can also be calculated numerically, but the measurement must
first be given a spatial implementation.  For $n_{\rm shell}=1$, every
controller orbital has support on at most $3\times3$ cells and two microscopic
orbitals per cell.  Choose a fixed nearest-neighbor complex-Givens tree
$W_e(\phi)$ that rotates $|\chi_e(\phi)\rangle$ onto an anchor orbital at the
controller center.  One exact event is then compiled as
\begin{equation}
  C_e\xrightarrow{W_e(\phi_e)}C_{e,1}
  \xrightarrow{\text{local number measurement and fSWAP reset}}C_{e,2}
  \xrightarrow{W_e(\phi_e)^\dagger}C_{e+1}.
  \label{eq:local-measurement-compilation}
\end{equation}
The local readout and its collocated fresh ancilla exchange charge with the
reservoir but do not define transverse motion.  Spatial transport occurs
during the explicitly routed charge-conserving Givens gates.

For every elementary one-particle gate $u_g$ and half-system cut projector
$R$, accumulate
\begin{equation}
  \frac{q_{x,g}}{e}
  =\Tr\!\left[
    R\left(u_gC_gu_g^\dagger-C_g\right)
  \right].
  \label{eq:gate-resolved-current}
\end{equation}
Only gates whose declared spatial path crosses the cut contribute.  The
trajectory current is
\begin{equation}
  Q^{\rm gate}_{x,\bxi}=\sum_{g\in\bxi}q_{x,g}.
  \label{eq:compiled-trajectory-current}
\end{equation}
Unlike the reduced balance Eq.~\eqref{eq:microscopic-transverse-transfer},
Eq.~\eqref{eq:compiled-trajectory-current} does not count a Bayesian
measurement innovation as transported charge.  The same gate contractions on
the deterministic covariance Eq.~\eqref{eq:exact-born-mean-replacement} give
the exact Born-mean electrical current.

The compilation is part of the microscopic model.  Two local Givens trees can
realize the same reduced rank-one update while distributing short-distance
currents differently.  A quantized large-scale coefficient should agree
between reasonable local trees and under cut deformation.  Required
validation is eventwise equality with the current rank-one engine at zero
field---identical outcome probability and final covariance to numerical
tolerance---together with gate-resolved continuity on the enlarged
system--ancilla circuit.  The reduced source-subtracted pump is the inexpensive
discovery calculation; the compiled gate current is the decisive calculation
for the phrase ``electrical Hall conductance'' and for a strong
record-concentration claim.

\subsection{Immediately runnable sampled finite-step estimator}

The current engine can already measure a microscopic linear impulse response
without any dynamics modification.  First draw a Born batch at zero twist and
burn it in to $G_{0,\bxi}$.  From every same baseline covariance, launch two
ordinary continuations with static seam twists $+\epsilon$ and $-\epsilon$.
The outcomes in both continuations are newly Born sampled under their own
twisted covariances; they must never be frozen to the zero-twist word.

A sudden twist step is an impulse of the circumferential electric field.  With
an orientation sign $\varsigma=\pm1$ fixed once from the coordinate and flux
conventions, the direct central estimator is
\begin{equation}
  \widehat C_{\rm step}(c;\epsilon)
  =\varsigma\frac{\pi}{\epsilon}
  \left[
    \overline{q_x^{(+\epsilon)}(c)}
    -\overline{q_x^{(-\epsilon)}(c)}
  \right].
  \label{eq:microscopic-step-estimator}
\end{equation}
The sign $\varsigma$ is not chosen after seeing the data.  Reversing the twist
must reverse the unscaled transfer, while reversing the programmed Chern
orientation must reverse the inferred integer.  A linear-response result
requires an $\epsilon^2$ extrapolation of the central estimate,
\begin{equation}
  \widehat C_{\rm step}(\epsilon)
  =C_0+b\epsilon^2+O(\epsilon^4),
  \label{eq:step-extrapolation}
\end{equation}
and a time window in which the integrated response has converged before finite
size recurrence or uncontrolled late-time relaxation.

The executable sequence is:
\begin{enumerate}
  \item run \code{run\_markov\_circuit} at $\phi=0$ with
  \code{G\_history=False}, \code{save=False}, and
  \code{return\_data=True}; retain only \code{G\_final};
  \item snapshot the GPU and NumPy random states immediately before the
  continuation;
  \item build the $+\epsilon$ model, call
  \code{set\_controller\_twist}, restore the snapshot, and continue from
  \code{G\_init=G\_final} with random schedules and ordinary Born outcomes;
  \item repeat from the identical baseline and random-state snapshot at
  $-\epsilon$;
  \item use the cycle observer for Eq.~\eqref{eq:regional-particle-number} and
  the record observer with a precomputed table
  Eq.~\eqref{eq:regional-controller-weight} for
  Eq.~\eqref{eq:microscopic-ancilla-source};
  \item construct Eq.~\eqref{eq:microscopic-step-estimator} separately for
  every independent parent trajectory, then average and bootstrap whole
  trajectory blocks.
\end{enumerate}
Restoring the same random stream is a common-random-number variance reduction,
not postselection.  Each marginal run still compares its shared uniform random
number with its own Born probability.  Branches may and should diverge between
$+\epsilon$ and $-\epsilon$.  Perfect correction has a fixed number of random
draws per event, making this coupling particularly clean.

A sensible first grid is
\begin{equation}
  \epsilon\in\{\pi/8,\pi/16,\pi/32\},
  \qquad c\le 2N_y,
  \label{eq:step-pilot-grid}
\end{equation}
rather than $\epsilon=10^{-3}$.  For $C=1$, the raw one-sided signal is only
$|q_x|\simeq\epsilon/(2\pi)$; an unnecessarily tiny step converts ordinary
trajectory noise into an impossible sample requirement.  The three finite
steps allow the linear limit to be tested rather than assumed.

\subsection{Direct quantized pump: the stronger numerical experiment}

The clearest quantization test avoids dividing a small noisy signal by
$\epsilon$.  After zero-twist burn-in, ramp the microscopic controller
holonomy through one full circle,
\begin{equation}
  \phi_c=2\pi f(c/T_\phi),
  \qquad
  f(u)=u-\frac{\sin(2\pi u)}{2\pi},
  \qquad c=0,\ldots,T_\phi,
  \label{eq:smooth-flux-ramp}
\end{equation}
where the smooth profile has zero endpoint velocity.  At the beginning of
each cycle, twist the controller orbitals to $\phi_c$, visit a newly random
site schedule, sample the actual Born outcomes, perform feedback, and update
Eqs.~\eqref{eq:microscopic-ancilla-source}--
\eqref{eq:microscopic-transverse-transfer}.  The provisional reduced-channel
per-record pump estimator is
\begin{equation}
  \widehat C_{{\rm pump},\bxi}
  =\varsigma q_{x,\bxi}(T_\phi),
  \label{eq:microscopic-pump-estimator}
\end{equation}
because a $2\pi$ cycle transfers $C$ particles in the ideal orientation.
This is an order-one signal and is the recommended discovery observable.  For
the final electrical claim, replace $q_x$ by
$Q^{\rm gate}_{x,\bxi}/e$ from
Eq.~\eqref{eq:compiled-trajectory-current}.

The current static-twist method already supplies all projector phases.  A
zero-engine-change pilot can run the entire ramp in one uninterrupted call to
\code{run\_markov\_circuit}: the cycle observer records the state at cycle
$c$ and then calls
\code{set\_controller\_twist(phi[c+1], gauge="seam")} so the next cycle uses
the next absolute flux.  The initial $c=0$ observer call installs the first
ramp value.  This preserves the canonical cycle loop, newly random schedules,
and live Born draws.  It also avoids the initialization overhead and semantic
ambiguity of repeated one-cycle runs.

For production provenance, the minimal canonical-engine extension is still a
validated \code{controller\_twist\_schedule} argument, applied at the start of
each cycle before its measurement updates.  A seam-gauge ramp is the primary
protocol.  A time-dependent uniform-gauge comparison must not be advertised
until the associated temporal gauge/frame kick has been implemented.  In seam
gauge, the weights $\Tr(R_aP_e)$ are independent of $\phi$ because the twist
changes component phases but not their magnitudes; they can be precomputed
once.  H3 cannot be repurposed: H3 rebuilds a static model at each flux and
forces one zero-flux record, whereas Eq.~\eqref{eq:smooth-flux-ramp} is one
continuous physical trajectory with fresh Born outcomes.

Run both orientations of the flux cycle and form the flux-odd estimator
\begin{equation}
  \widehat C^{\rm odd}_{\rm pump}(T_\phi)
  =\frac{\varsigma}{2}
  \left[
    \overline{q_x^{(+2\pi)}(T_\phi)}
    -\overline{q_x^{(-2\pi)}(T_\phi)}
  \right].
  \label{eq:microscopic-odd-pump-estimator}
\end{equation}
It cancels flux-even stochastic drift and is the primary reduced-channel
plateau estimator.  The individual forward and reverse distributions remain
mandatory outputs.

For each $N_y$ and $T_\phi$, report
\begin{align}
  \bar C_{\rm pump}
  &=\frac1S\sum_{a=1}^{S}\widehat C_{{\rm pump},a},\nonumber\\
  V_C
  &=\frac1{S-1}\sum_{a=1}^{S}
    (\widehat C_{{\rm pump},a}-\bar C_{\rm pump})^2,\nonumber\\
  f_{\rm bad}(\delta)
  &=\frac1S\sum_{a=1}^{S}
    \mathbf{1}\!\left(
      |\widehat C_{{\rm pump},a}-C_{\rm target}|>\delta
    \right).
  \label{eq:microscopic-pump-statistics}
\end{align}
The indicator in Eq.~\eqref{eq:microscopic-pump-statistics} is an ordinary
classical indicator, not an identity operator.  The Born sample mean tests a
no-postselection plateau; decreasing $V_C$ and $f_{\rm bad}$ with size and
ramp time test record typicality.

\subsection{Concrete pilot and production ladders}

The first calculation is an exact Born-mean screen using
\code{mean\_replacement=True}, with a fixed random schedule per batch member.
It requires only a lightweight event observer that emits
Eq.~\eqref{eq:exact-mean-transfer}; the flux itself can be advanced by the
cycle callback just described.  If this deterministic mean has no plateau,
large trajectory sampling is not justified.

The first sampled no-engine-change pilot is:
\begin{itemize}
  \item explicit interface, $N_x\times N_y=20\times24$,
  $n_{\rm shell}=1$, perfect correction, random schedule;
  \item $S=10$ independent baseline trajectories and $2N_y$ burn-in cycles;
  \item paired $\pm\epsilon$ continuations for
  $\epsilon=\pi/8,\pi/16$, each observed for $2N_y$ cycles;
  \item the matched-trivial controller and reversed twist as mandatory
  controls.
\end{itemize}
This costs one burn-in plus four continuations, approximately five ordinary
$2N_y$ parent runs, but they remain sample-batched on the A100.  It stores only
regional charge, source, continuity residuals, and compact record diagnostics;
no covariance histories are needed.

If the signal and signs pass, add $\epsilon=\pi/32$, use $S=25$, and repeat at
$N_y=24,32$ before the full size ladder.  The ramp discovery run should start
at $20\times24$, $S=5$, and $T_\phi=4N_y,8N_y$, followed by a pilot with
$S=10$ and $T_\phi=4N_y,8N_y,16N_y$.  Production uses $S=25$ in shards of
five, the interface forward ramp at all three rates, the interface reverse
ramp at $8N_y$, and matched-trivial forward/reverse controls at $8N_y$.  A
$20\times32$ size check is promoted only after the measured timing and
variance are acceptable.  The required extrapolations are
\begin{equation}
  T_\phi^{-1}\to0,
  \qquad N_y^{-1}\to0,
  \qquad S^{-1/2}\to0\ \text{statistically},
  \label{eq:pump-limits}
\end{equation}
with burn-in convergence checked separately.  A finite set of 25 records can
support a plateau and an initial distributional statement, not a theorem about
every record.

The numerical acceptance checks are:
\begin{enumerate}
  \item $q_L+q_R=0$ and
  $\Delta N_{\rm total}=\sum_e(s_e-m_e)$ within floating-point tolerance;
  \item odd response under flux reversal and Chern-orientation reversal;
  \item a matched-trivial result converging to zero;
  \item stability under moving both the twist seam and the two regional cuts;
  \item agreement of seam and properly transformed uniform gauges;
  \item convergence in step size or ramp rate, circumference, burn-in, and
  sample count;
  \item retention of the complete per-record distribution rather than only
  its mean.
\end{enumerate}

This microscopic program is sufficient by itself.  The generating-functional
language below is retained only as an optional analytic consistency check and
is not part of the proposed numerical implementation.

\subsection{Why the microscopic derivative is legitimate outside equilibrium}

Linear response is a functional derivative around a reference dynamics, not
an equilibrium privilege.  Equilibrium supplies the Kubo--Martin--Schwinger
condition, fluctuation--dissipation relations, and familiar spectral proofs of
quantization.  A causal response around a transient, periodic state, or
nonequilibrium stationary path ensemble remains well defined without any of
those simplifications
\cite{VenutiZanardi2016,KonopikLutz2019,TalkingtonClaassen2024}.

For one explicit space--time realization the kernel
\begin{equation}
  K^R_{xy}(n,n')
  =\left.
  \frac{\delta\langle I_x(n)\rangle}{\delta A_y(n')}
  \right|_{A=0},
  \qquad n>n',
  \label{eq:two-time-hall-kernel}
\end{equation}
depends on both times separately.  After burn-in, stationarity of the Born path
measure permits an ensemble Fourier transform in $n-n'$.  Statistical
stationarity is sufficient; an individual measurement word need not repeat.
The random schedule and intrinsic Born outcomes therefore act as genuine
exogenous and endogenous space--time disorder, respectively.

\subsection{Optional generating-functional check}

Let $K_{\bxi}[A]$ be the complete history-dependent branch operator, including
feedback conditioned on earlier entries of $\bxi$.  Couple the forward and
backward branches to independent gauge fields $A_+$ and $A_-$, and insert a
counting field $\chi_x$ into every elementary operation that transports charge
through a declared $x$ cut.  The exact instrument generating functional is
\begin{equation}
  \mathcal Z[A_+,A_-;\chi_x]
  =\sum_{\bxi}q(\bxi_{\rm ext})
  \Tr\!\left[
    K_{\bxi}[A_+;+\chi_x/2]\rho_{\ssub}
    K_{\bxi}[A_-;-\chi_x/2]^\dagger
  \right],
  \label{eq:instrument-keldysh-functional}
\end{equation}
where $q(\bxi_{\rm ext})$ is the field-independent probability of any
exogenous schedule choices.  The same classical record labels join the two
contour branches, while all intrinsic outcomes are summed.  Trace preservation
gives the unitarity identity
\begin{equation}
  \mathcal Z[A,A;0]=1.
  \label{eq:keldysh-normalization}
\end{equation}
Thus Eq.~\eqref{eq:instrument-keldysh-functional} is a literal discrete-time
Schwinger--Keldysh construction for the adaptive instrument.  It is not a
postselected or representative-record object.  Replica Keldysh formulations
for monitored jump processes and hybrid classical--quantum full counting
statistics provide closely related precedents
\cite{KloiberTollingerSieberer2026,RosalFiusaPottsLandi2026}.

Define
\begin{equation}
  A_c=\frac{A_++A_-}{2},
  \qquad A_q=A_+-A_-,
  \qquad W=-\ii\log\mathcal Z.
\end{equation}
Up to the declared sign and charge normalization, the physical current and
retarded response are
\begin{align}
  \bar j_x(n)
  &=\left.\frac{\delta W}{\delta A_{x,q}(n)}\right|_{A_q=0},
  \label{eq:keldysh-current}\\
  K^R_{xy}(n,n')
  &=\left.
  \frac{\delta^2W}
  {\delta A_{x,q}(n)\delta A_{y,c}(n')}
  \right|_{A_q=0}.
  \label{eq:keldysh-retarded-response}
\end{align}
If $A_y$ enters only through a Hamiltonian term $-A_yJ_y$, this reduces to a
retarded current--current correlator plus the diamagnetic/contact term.  In the
present circuit the field also changes controller orbitals, projectors, Kraus
maps, and feedback operations.  The exact insertion is therefore an
instrument vertex $\partial_{A_y}\cE$, not generically a second current
operator.  A symmetrized current autocorrelation measures noise; outside
equilibrium no fluctuation--dissipation theorem converts it into
Eq.~\eqref{eq:keldysh-retarded-response}.

\subsection{The Born score inside the Keldysh functional}

At equal physical fields,
\begin{equation}
  \bar I_x[A]=\sum_{\bxi}p_A(\bxi)I_{x,\bxi}[A].
\end{equation}
Differentiating gives
\begin{equation}
  \left.\frac{\delta\bar I_x}{\delta A_y}\right|_0
  =\E_0\!\left[ K^{\cond}_{xy,\bxi}\right]
   +\E_0\!\left[
     (I_{x,\bxi}-\bar I_x)S_{y,\bxi}
   \right],
  \label{eq:hall-score-decomposition}
\end{equation}
with
\begin{align}
  K^{\cond}_{xy,\bxi}
  &=D_GI_x[X_y]+\partial_{A_y}I_x,
  \label{eq:current-response}\\
  S_{y,\bxi}
  &=\left.\delta_{A_y}\log p_A(\bxi)\right|_0.
\end{align}
The gauge tangent has the inhomogeneous recursion
\begin{equation}
  X_{y,e+1}=D_G\Phi_e[X_{y,e}]+\partial_{A_y}\Phi_e.
  \label{eq:gauge-tangent}
\end{equation}
When the controller projector also depends on $A_y$, both the map tangent and
the score contain the corresponding $\dot P$ terms.  Equation
\eqref{eq:hall-score-decomposition} is exactly what the doubled functional
produces: the first term moves the normalized state along a fixed word, while
the second differentiates that word's Born weight.  The score is therefore not
an optional correction.  It is the classical record-sector part of the
physical no-postselection response.

For Monte Carlo analysis one can retain the single-record influence estimator
\begin{equation}
  \widehat K_{xy,\bxi}
  =K^{\cond}_{xy,\bxi}
   +(I_{x,\bxi}-\bar I_x)S_{y,\bxi},
  \qquad
  \E[\widehat K_{xy,\bxi}]=K^{\Born}_{xy}.
  \label{eq:hall-influence-estimator}
\end{equation}
Its distribution is meaningful for diagnosing typicality, but its score term
is an influence-function contribution, not the response of a postselected
branch.

\subsection{The model-specific conservation obstruction}

The reduced physical fermion system is not charge conserving.  When a
measurement outcome disagrees with its controller target, perfect correction
uses a fresh filled or empty ancilla and an fSWAP.  Total $U(1)$ charge is
conserved only in the enlarged system--ancilla implementation.  The physical
system obeys a sourced circuit-time continuity equation,
\begin{equation}
  \Delta n_{\br}+(\nabla_{\rm lat}\!\cdot j)_{\br}
  =s^{\rm anc}_{\br},
  \label{eq:sourced-continuity}
\end{equation}
not a source-free one.  Consequently the current reduced covariance engine and
its controller twist do not, by themselves, define a conventional
electromagnetic Hall current.

There are two honest observables:
\begin{enumerate}
  \item \emph{Electrical Hall response of a specified dilation.}  Assign
  physical paths to every basis-rotation gate and ancilla port, Peierls-dress
  every charge-moving gate and fSWAP, and insert the counting field across all
  paths crossing the cut.  Equation~\eqref{eq:sourced-continuity} then becomes
  source free on the enlarged space.
  \item \emph{Generalized monitored-circuit pump.}  Use the controller
  holonomy as a parameter and explicitly subtract the known ancilla source.
  This is immediately useful, but should not be called electrical Hall
  conductance until its electromagnetic dilation is fixed.
\end{enumerate}
This distinction is physical: inequivalent local Stinespring compilations can
realize the same reduced Kraus map while routing microscopic charge along
different paths.  An electrical current is therefore not an intrinsic property
of the reduced instrument alone.  Independence of the quantized coefficient
from reasonable local compilations would itself be a nontrivial robustness
test.

At perfect correction the source subtraction is especially simple.  If the
measured occupation is $m\in\{0,1\}$ and the reset target is $s\in\{0,1\}$,
then
\begin{equation}
  C_{\rm after}-C_{\rm meas}=(s-m)P.
\end{equation}
For a one-particle region projector $R$, the exact injected charge is
\begin{equation}
  S^{\rm anc}_{R,e}=e(s_e-m_e)\Tr(RP_e).
  \label{eq:regional-ancilla-source}
\end{equation}
The executable already records $s_e-m_e$ as the event transfer.  With inward
current positive, a practical source-subtracted regional transfer is
\begin{equation}
  Q^{\rm sub}_{\partial R}
  =e\,\Delta N_R-\sum_{e\in\mathrm{ramp}}S^{\rm anc}_{R,e}.
  \label{eq:source-subtracted-current}
\end{equation}
Equation~\eqref{eq:regional-ancilla-source} is exact for perfect correction.
If feedback can fail, the code must archive the actual reset transfer rather
than merely the requested target minus outcome.

\subsection{What quantization would mean in Keldysh language}

In units $e=\hbar=1$, the desired parity-odd long-distance functional is
\begin{align}
  W_{\rm top}
  &=\frac{C}{4\pi}\int\!\dd^3x\,
    \epsilon^{\mu\nu\lambda}
    \left(
      A_{\mu,q}\partial_\nu A_{\lambda,c}
      +A_{\mu,c}\partial_\nu A_{\lambda,q}
    \right)\nonumber\\
  &=\frac{C}{2\pi}\int\!\dd^3x\,
    \epsilon^{\mu\nu\lambda}
    A_{\mu,q}\partial_\nu A_{\lambda,c}
    +\text{boundary term}.
  \label{eq:keldysh-chern-simons}
\end{align}
It yields
\begin{equation}
  j_x=\frac{C}{2\pi}E_y,
  \qquad \sigma_{xy}=C\frac{e^2}{h}.
  \label{eq:quantized-hall-response}
\end{equation}
Additional causal dissipative kernels and terms quadratic in $A_q$ encode
relaxation and noise.  They do not by themselves forbid the topological term.
Keldysh Chern--Simons response has been established for specially structured
number-conserving open systems with a pure stationary state
\cite{TonielliBudichAltlandDiehl2020}.  That result is a template, not a theorem
for the present circuit: pure conditioned trajectories generally average to a
mixed Born state unless all positive-probability branches lie on the same ray.

For this adaptive process integer $C$ is therefore a conjectural response
coefficient that must be demonstrated under explicit hypotheses:
\begin{enumerate}
  \item exact $U(1)$ gauge covariance of the full local dilation;
  \item locality or controlled quasilocality of all charge-moving operations;
  \item a stationary, ergodic late-time Born process;
  \item a bulk dynamical, mobility, or rapid-mixing gap that remains controlled
  throughout the twist cycle;
  \item vanishing longitudinal DC response in the scaling limit;
  \item large-gauge closure and invariance under seam and cut deformation;
  \item a stated order of transient, field, slow-ramp, thermodynamic, and
  statistical limits.
\end{enumerate}
Static disorder can coexist with almost-sure quantized Hall response under
mobility-gap hypotheses
\cite{NiuThoulessWu1985,BellissardVanElstSchulzBaldes1994}.  Quantized Floquet
pumping can survive spatial disorder and, over a controlled regime, temporal
noise
\cite{TitumBergRudnerRefaelLindner2016,TimmsSiebererKolodrubetz2021}.
These are important precedents, but their disorder is not the endogenous,
state-dependent Born record of an adaptive instrument.

The caution is equally important.  A conserved state Chern number after a
quench need not imply quantized Hall response, and dissipatively prepared
topology and transport can separate
\cite{CaioCooperBhaseen2016,ShavitGoldstein2020}.  Thus the proposed experiment
would not merely remeasure the known trajectory invariant.

\subsection{Three meanings of record-independent quantization}

The following claims have strictly increasing strength:
\begin{enumerate}
  \item \emph{Born-mean plateau:}
  \begin{equation}
    \frac{\E_{\Born}[Q_{x,\bxi}]}{e}\longrightarrow C.
    \label{eq:born-mean-pump}
  \end{equation}
  This is already a physical no-postselection result, but it can hide a broad
  or bimodal record distribution.
  \item \emph{Born-typical concentration:}
  \begin{equation}
    \Pr_{\Born}\!\left(
      \left|Q_{x,\bxi}/e-C\right|>\delta
    \right)\longrightarrow0
    \quad\text{for every fixed }\delta>0.
    \label{eq:record-concentration}
  \end{equation}
  This is the natural analogue of an almost-sure disorder statement and the
  strongest realistic meaning of ``record independent.''
  \item \emph{Uniform every-word quantization:} every allowed finite record,
  including arbitrarily rare words, gives exactly the same integer.  This
  would require a uniform gap and is neither needed nor currently plausible.
\end{enumerate}
A quantized linear-response coefficient does not logically require the entire
finite-time counting distribution to collapse to a delta function; irreducible
shot noise can remain.  The full counting statistics should therefore report
the mean response, response variance, higher cumulants, quantiles, and
bad-record fraction separately.  Exact claims concern specified asymptotic
limits, never a finite sample of 25 trajectories.

\subsection{Dynamical bulk--boundary correspondence}

Suppose two bulk regions have Born Keldysh coefficients $C_L$ and $C_R$.  At
their interface, the gauge variation of Eq.~\eqref{eq:keldysh-chern-simons}
has an uncancelled boundary term proportional to
\begin{equation}
  \Delta C=C_L-C_R.
\end{equation}
Gauge invariance of the full dilation then requires a wall anomaly with net
chirality $\Delta C$.  This is the anomaly-inflow statement of
bulk--boundary correspondence, formulated for the Born-summed dynamical
functional rather than a ground-state Hamiltonian.

H2 tests one side of this relation: wall-localized, signed propagation.  It
does not measure the coefficient $C$.  H3 tests branchwise auxiliary
entanglement spectral flow under a frozen record.  It is not a physical flux
ramp and does not sum the flux-dependent Born process.  A quantized bulk pump
or Kubo coefficient agreeing with the H2 wall chirality would supply the
missing side.  If Eq.~\eqref{eq:record-concentration} also holds, the result
would be genuinely stronger: dynamical bulk--boundary correspondence would
survive endogenous space--time measurement disorder without postselection.

\subsection{Two operational protocols}

\subsubsection{Protocol A: full-dilation retarded Hall kernel}

Use a homogeneous topological torus or cylinder first, rather than the
two-wall geometry.  Burn in to the stationary Born process, apply a weak
time-dependent $A_y$ to every affected operation of the local dilation, and
count the full charge crossing an $x$ cut.  Evaluate
Eq.~\eqref{eq:keldysh-retarded-response} either from the exact inhomogeneous
state tangent plus score or from common-schedule $\pm A$ replays as a
validation.  After testing stationarity, extract the antisymmetric response,
take the thermodynamic or mobility limit before the DC limit, and verify
$\sigma_{xx}\to0$.

For a stationary augmented channel $\mathbb E_0$, the discrete causal kernel
has the transparent form
\begin{equation}
  K^R_{O,A}(\ell)
  =\Theta(\ell)
  \left\langle\!\left\langle O\middle|
    \mathbb E_0^{\ell-1}(\partial_A\mathbb E)_0
    \middle|\varrho_{\ssub}\right\rangle\!\right\rangle
  +K_{\rm contact}.
  \label{eq:augmented-channel-kubo}
\end{equation}
Classical feedback memory can be included in the augmented state.  This is the
discrete NESS analogue of a Lindblad resolvent, not an equilibrium fluctuation
formula.

\subsubsection{Protocol B: no-postselection Laughlin pump}

The numerically cleaner quantization test is a physical $y$-flux ramp
\begin{equation}
  \phi_y:0\longrightarrow2\pi,
  \qquad A_y=\phi_y/N_y.
\end{equation}
Outcomes must be newly sampled from the instantaneous flux-dependent Born
process; freezing them would convert the experiment back into H3.  Count the
full-dilation charge, or provisionally the source-subtracted charge
Eq.~\eqref{eq:source-subtracted-current}, crossing an $x$ cut.  With orientation
sign fixed in advance, the target is
\begin{equation}
  Q_x/e\longrightarrow\pm C.
  \label{eq:laughlin-pump-target}
\end{equation}
The associated full counting statistics can be generated by the tilted
instrument
\begin{equation}
  \cE_{A,\chi}(\rho)
  =\sum_m K_m[A;+\chi/2]\rho K_m[A;-\chi/2]^\dagger,
  \label{eq:tilted-instrument}
\end{equation}
with
\begin{align}
  \langle Q_x\rangle
  &=\left.\partial_{\ii\chi}\log\mathcal Z\right|_{\chi=0},\nonumber\\
  \partial_A\langle Q_x\rangle
  &=\left.\partial_A\partial_{\ii\chi}
    \log\mathcal Z\right|_{A=\chi=0}.
  \label{eq:fcs-derivatives}
\end{align}
Tilted generators and trajectory large deviations are established tools
\cite{GarrahanLesanovsky2010,FlindtNovotnyBraggioSassettiJauho2008}.

A discovery pilot can use the current $20\times24$ or $20\times32$ explicit
interface, $n_{\rm shell}=1$, perfect correction, random schedules, five to ten
Born trajectories, and ramp times $T_\phi=4N_y,8N_y,16N_y$ after a burn-in of
order $2N_y$.  It should stream wall/strip charge, total system charge, and
spatially resolved ancilla source rather than covariance histories.  This
pilot can establish feasibility and sign conventions; it cannot establish a
quantized thermodynamic coefficient.

\subsection{Controls and implementation consequences}

The minimum controls are:
\begin{enumerate}
  \item trivial controller, reversed Chern orientation, and reversed ramp;
  \item seam versus uniform gauge, seam translation, and cut deformation;
  \item a discrete Ward identity including the ancilla source, together with
  Eq.~\eqref{eq:keldysh-normalization} and $\E[S]=0$;
  \item random versus raster schedules and controlled measurement noise;
  \item probe-amplitude, frequency or ramp-rate, size, and wall-separation
  scaling;
  \item longitudinal response and full counting noise;
  \item at least two local physical compilations of the dilation if the result
  is called electrical Hall conductance.
\end{enumerate}

The existing static method \code{set\_controller\_twist} is useful but
insufficient.  A new versioned campaign would need a time-dependent twist
schedule inside the canonical GPU engine, ordinary Born sampling at each
instantaneous twist, a cut/source observer, and the $\dot P$ terms in the gauge
tangent and score.  A time-dependent uniform gauge also requires the proper
inter-cycle gauge or scalar-potential transformation; repeatedly rebuilding
static H3 models is not a physical ramp.  The implementation should validate
eventwise gauge derivatives against $\pm\delta\phi$, $2\pi$ closure,
seam/uniform equivalence, source-subtracted continuity, and exact small-system
branch enumeration before production.

Quantized response need not always be framed as conventional DC transport.
Frequency-integrated circular dichroism and certain pulse-to-density protocols
can also tie responses to topological invariants in closed systems
\cite{RepellinGoldman2019,Zhang2023}.  Those alternatives do not make the H2
density ridge quantized, nor do they remove the dilation issue for electrical
transport.

\section{Relation to the literature and novelty boundary}
\label{sec:literature}

Table~\ref{tab:literature} positions the construction without claiming that
monitored-system response theory itself is new.

\begin{table}[t]
\caption{Literature map and project-specific distinction.}
\label{tab:literature}
\small
\begin{tabularx}{\textwidth}{@{}L{0.21\textwidth}L{0.32\textwidth}Y@{}}
\toprule
Established framework & What is established & What remains specific here \\
\midrule
Closed Kubo and chiral edge theory
\cite{Kubo1957,Wen1990} & Retarded density response and a one-way pole
$\omega=vk$ for a chiral edge. & Adaptive normalized backaction, record
dependence, and the Born score. \\
Lindblad/NESS response
\cite{AvronFraasGraf2012,VenutiZanardi2016,AlbertBradlynFraasJiang2016,KonopikLutz2019,LevyRabaniLimmer2021}
& Generator linearization, resolvents, geometric response, and conditions under
which open Hall response can remain quantized. & Discrete rank-losing
record-conditioned Gaussian instrument and branch distribution. \\
Keldysh topology and monitored field theory
\cite{TonielliBudichAltlandDiehl2020,TalkingtonClaassen2024,KloiberTollingerSieberer2026}
& Nonequilibrium response functionals, Chern--Simons response in structured
open systems, and Keldysh representations of monitored processes. & A doubled,
history-dependent feedback instrument whose Born score and physical dilation
are kept explicit. \\
Gaussian process response \cite{MehboudiParrondoAcin2019} & Covariance-level
response of Gaussian channels. & Fermionic projective normalized map with
exact outcome conditioning and feedback. \\
Quantum trajectory likelihood and response
\cite{GammelmarkMolmer2013,Vu2025} & Record likelihoods, tangent inference,
Fisher information, and bounds on trajectory-observable response. & Local
density impulse, exact branch congruence, and score-completed wall
susceptibility in this circuit. \\
Trajectory large deviations
\cite{GarrahanLesanovsky2010} & Tilted generators and statistics of record
observables. & Spatially resolved class-A domain-wall response and its
conditional/Born comparison. \\
Feedback and monitored topology
\cite{NakagawaUeda2025,OshimaMochizukiHamazakiFuji2025,XiaoKawabata2026} &
Topological quantum channels, monitored Lyapunov spectra, and boundary modes. &
Density-response spectroscopy of finite-range adaptive Gaussian Born
trajectories, including rank loss and score reweighting. \\
Dissipative topology
\cite{Goldstein2019,ShavitGoldstein2020,GneitingKoottandavidaRozhkovNori2022}
& Constraints on dissipative preparation and examples of trajectory/open-system
topology. & Direct comparison among fixed-record, Born, mean-channel, and
frozen-holonomy diagnostics on the same programmed wall. \\
Disordered topological transport
\cite{NiuThoulessWu1985,BellissardVanElstSchulzBaldes1994,TitumBergRudnerRefaelLindner2016,TimmsSiebererKolodrubetz2021}
& Almost-sure static-disorder quantization and dynamical pumping robust to
spatial or temporal disorder under controlled hypotheses. & Endogenous Born
space--time disorder, score-completed response, and record-concentration tests
in an adaptive circuit. \\
\bottomrule
\end{tabularx}
\end{table}

The plausible project-specific contribution is therefore a synthesis:
\begin{quote}
An exact, record-conditioned tangent response for a normalized, rank-losing
adaptive fermionic Gaussian instrument, completed by the many-body Born
likelihood score and analyzed as a wall-resolved momentum--frequency density
susceptibility across a topological trajectory ensemble.
\end{quote}
The more ambitious follow-on contribution would be a separate statement:
\begin{quote}
A gauge-covariant, no-postselection adaptive fermion instrument has a Born Hall
coefficient locked to an integer, with its response distribution concentrating
over endogenous space--time measurement records and matching the chiral wall
anomaly.
\end{quote}
Nothing in the present trajectory Chern data, H2 ridge, or H3 frozen holonomy
already proves this second statement.
The individual ingredients are not new.  A formal novelty claim should use
``to our knowledge'' only after a broader systematic search.  In particular,
recent topology of discrete quantum feedback
\cite{NakagawaUeda2025} is a close neighboring field that must be discussed,
while monitored free-fermion classifications
\cite{XiaoKawabata2026} supply the appropriate tangent/Lyapunov comparison.

\section{Current implementation and exact-tangent upgrade}
\label{sec:implementation}

\subsection{Current production contract}

The current H2 production choices are listed in
Table~\ref{tab:h2-contract}.  They are taken from the audited working campaign
and executable configuration, not from the superseded numerical-campaign copy
inside the draft directory.

\begin{table}[t]
\caption{Current H2/H3 contract as of August 18, 2026.}
\label{tab:h2-contract}
\small
\begin{tabularx}{\textwidth}{@{}L{0.24\textwidth}Y@{}}
\toprule
Quantity & Current choice \\
\midrule
H2 geometry & $N_x=20$; $\alpha_{\rm top}=1$,
$\alpha_{\rm triv}=30$; $n_{\rm shell}=1$; topological columns
$x=5,\ldots,15$; interfaces on bonds $(4,5)$ and $(15,16)$; periodic $y$ \\
H2 circumference & $N_y=32,48,64$ \\
Independent trajectories & $S=25$, in five shards of five \\
Schedule law & Random serial only \\
Physical duration & $2N_y$ cycles from independent product-state starts \\
Sources & Eight translated $y_0$ positions on each of two wall orientations \\
Time origins & Four late origins per parent trajectory \\
Response horizon & $N_y/2$ \\
Wall window & Half-width one cell in current executable config \\
Saved decomposition & Conditional derivative, likelihood-score contribution,
total Born response, and within-trajectory source variance.  The cumulative
scalar score itself is not currently archived. \\
Current derivative & Central $\pm\epsilon$ replay with
$\epsilon=10^{-3}$; validation at $5\times10^{-4}$ and $2\times10^{-3}$ \\
Statistics & Whole parent trajectory is the bootstrap unit \\
H3 & One interface and one matched-trivial record at $20\times24$; 65 twist
points per closed circle; seam $y=23\to0$; half-torus entanglement spectrum with
wall-weight classification \\
\bottomrule
\end{tabularx}
\end{table}

The H2 runner already preserves the essential estimator order: it freezes each
baseline schedule and outcome record, central differences the conditional
density and its chronological log likelihood, multiplies the latter by the
baseline density, translates sources within each trajectory, and retains the
trajectory axis for later bootstrapping.  It does not yet propagate
Eqs.~\eqref{eq:tangent-recursion}--\eqref{eq:score-recursion}.  The repository
also does not yet contain the final downstream H2 Fourier/ridge merger; the
current compact product is the real-space input to that analysis.  A separate
geometry caveat is also recorded here: the canonical CPU class currently uses
an $N_x/3$ default wall width, whereas the production GPU H2/H3 geometry uses
$N_x/4$.  Production-response claims must use the GPU geometry stated in
Table~\ref{tab:h2-contract}, or explicitly reconcile the defaults before
combining CPU and GPU results.

\subsection{Recommended exact-tangent implementation}

The upgrade should preserve the immutable finite-difference campaign rather
than silently changing its schema.  A new versioned H2 product should:
\begin{enumerate}
  \item initialize a batched tangent $X_0=-\ii[P_0,G]$ for every live source;
  \item propagate the base covariance once through the canonical GPU engine;
  \item update all source tangents with the exact local or dense rank-one
  Jacobian inside the canonical GPU class, respecting repository policy;
  \item accumulate the scalar score increment
  Eq.~\eqref{eq:binary-score} on the same device;
  \item stream only $\Tr(P_rX)/2$, the baseline density, score, wall leakage,
  denominator minima, and compact diagnostics;
  \item archive an explicit schema/version and the exact engine hash;
  \item compare against the existing $\pm\epsilon$ implementation without
  overwriting either result.
\end{enumerate}
This removes two full continuations per source and should materially reduce H2
runtime, although the realized speedup depends on how many source tangents fit
in the A100 batch and must be benchmarked rather than assumed.

\subsection{Validation gates}

Before the analytic tangent replaces the finite difference as the primary
estimator, require:
\begin{enumerate}
  \item event-level agreement of Eq.~\eqref{eq:event-tangent} with central
  differences for occupied and empty outcomes, mixed and nearly pure states,
  local and dense backends, and gain and loss targets;
  \item complete frozen-record agreement over a convergent $\epsilon$ window;
  \item agreement of the accumulated score with central differences of the
  complete chronological log probability;
  \item exact small-system branch enumeration of
  Eq.~\eqref{eq:score-identity};
  \item $\E[S]=0$ within whole-trajectory uncertainty, with the centered score
  estimator used in production analysis;
  \item CPU/GPU parity for identical frozen words;
  \item explicit counts of zero-probability/censored branches and minimum Born
  denominators;
  \item sign-convention parity with Eq.~\eqref{eq:implemented-kick}.
\end{enumerate}

\section{Analysis and claim protocol}
\label{sec:claim-protocol}

\subsection{Per-trajectory construction before averaging}

For each independent parent trajectory and each wall:
\begin{enumerate}
  \item average translated sources and permitted late origins only as correlated
  variance reduction inside that trajectory;
  \item retain $\chi^{\cond}_{w,\bxi}(r,t)$,
  $\chi^{\mathrm{score}}_{w,\bxi}(r,t)$, and their sum;
  \item record transverse wall retention and bulk leakage;
  \item Fourier transform each component with identical conventions;
  \item extract ridge slope, linewidth, wedge asymmetry, and fit quality;
  \item only then form ensemble means and bootstrap intervals over complete
  parent trajectories.
\end{enumerate}
Spatial sites, sources, time origins, and the two walls are not independent
Monte Carlo samples.

\subsection{Primary figures}

A compact but complete H2 presentation contains:
\begin{enumerate}
  \item real-space $\chi_w(r,t)$ heatmaps for both wall orientations;
  \item the full Born $\chi_w(k,\omega)$ for both walls, with common axes;
  \item an inset decomposing conditional and score contributions;
  \item the distribution of $v_{w,\bxi}$ or $\mathcal A_{w,\bxi}$;
  \item a matched trivial or bulk-source control;
  \item a size panel for signed velocity, wrong-chirality weight, and wall
  retention.
\end{enumerate}

\subsection{Acceptance and downgrade language}

The strong H2 statement requires a low-energy ridge whose velocity sign is
stable, reverses between the paired walls, remains wall localized, survives
whole-trajectory uncertainty and size scaling, and is absent in the matched
control.  Both terms in Eq.~\eqref{eq:full-density-response} belong in the
headline Born response.

If only $\chi^{\fixed}$ is chiral, report a chiral fixed-record conditional
response and state that the physical Born response is unresolved or different.
If only the mean is chiral while per-record signs are broad, do not claim that
typical individual trajectories share the same response.  If real-space drift
is visible but the Fourier ridge is unresolved, report directional spreading,
not a measured dispersion.  If a ridge exists without wall reversal or a
control null, report schedule-induced directional response rather than
topological chirality.  At no point should the H2 ridge be called a quantized
Hall conductance.

\section{Coherent research questions exposed by the framework}
\label{sec:questions}

The discussion suggests several questions that are logically independent and
could support a later paper beyond the immediate H2 analysis.

\subsection{Conditional versus Born chirality}

How large is
\begin{equation}
  \Delta\chi=\chi^{\Born}-\chi^{\fixed}
  =\operatorname{Cov}(n,S),
\end{equation}
and can the score change the direction, velocity, or only the spectral weight?
The answer measures how strongly the local perturbation reshapes the record
distribution rather than merely changing conditional states.

\subsection{Self-averaging and trajectory typicality}

Does the distribution of $v_{w,\bxi}$ narrow with $N_y$ and observation time,
or does a broad record-dependent chiral response survive?  A narrow
distribution would connect the Born response to an almost-sure trajectory
statement; a broad distribution would make the mean insufficient as a
description of typical monitored runs.

\subsection{Response poles versus tangent Lyapunov modes}

Does the linewidth $\gamma_w(k)$ extracted from
$\chi(k,\omega)$ agree with decay rates in the wall-localized tangent cocycle?
The two use the same normalized derivative but different source and projection
protocols.  Agreement would connect observed chiral propagation to the slow
mechanism responsible for purification; disagreement would identify distinct
sectors.

\subsection{Trajectory, mean-channel, and feedback-channel topology}

Which wall features survive after discarding the record?  Monitored Lyapunov
topology, dissipative topology, and discrete feedback-channel topology are
active neighboring subjects
\cite{GneitingKoottandavidaRozhkovNori2022,NakagawaUeda2025,OshimaMochizukiHamazakiFuji2025,XiaoKawabata2026}.
The present circuit can compare all three descriptions on the same microscopic
controller instead of assuming their equivalence.

\subsection{Record Fisher information and response cost}

The same derivatives entering the score determine the Fisher information
carried by the record.  One may ask whether a robust wall response requires the
record to learn the injected perturbation, and whether trajectory-response
bounds \cite{Vu2025} become tight near the wall.  This is a question about the
classical information channel of monitoring, not a substitute for the density
susceptibility.

\subsection{Gauge response and quantized pumping}

Can a gauge-covariant extension of the adaptive instrument support a robust
integer Born pump, and how is it related to the representative fixed-record
H3 flow?  This requires the full program of Sec.~\ref{sec:quantization}.  It is
a plausible separate work, not a prerequisite for demonstrating chirality in
the current manuscript.

\section{Recommended paper-level statement}
\label{sec:recommendation}

If the planned data pass their controls, the clean central statement is:
\begin{quote}
A local number-phase impulse produces a wall-localized density response with a
single low-energy propagation direction.  The momentum--frequency ridge
reverses between the two oppositely oriented domain walls.  This directionality
is visible in the distribution of fixed-record trajectory responses and
survives the Born likelihood-score completion required for the physical
ensemble.
\end{quote}
This statement is stronger than static criticality and more operational than
modular-time motion, while remaining narrower and safer than a quantized
transport claim.  H1 modular drift and H3 frozen-record spectral flow can be
presented as complementary state-level and branch-topology evidence.  The
conditional, fixed-record averaged, Born, and unconditional objects should be
named explicitly every time; ``trajectory average'' without a definition is no
longer adequate.

\appendix

\section{Derivation of the fixed-projector covariance congruence}
\label{app:derivative}

Let
\begin{equation}
  A=\Id+\eta GP,
  \qquad \Phi(G)=QA^{-1}GQ+\eta_sP.
\end{equation}
At fixed $P$, $\delta A=\eta XP$ and
\begin{align}
  D\Phi[X]
  &=Q\left[-A^{-1}(\eta XP)A^{-1}G+A^{-1}X\right]Q\nonumber\\
  &=QA^{-1}X\left[\Id-\eta PA^{-1}G\right]Q.
  \label{eq:raw-derivative}
\end{align}
The rank-one identities
\begin{align}
  QA^{-1}&=Q-\frac{\eta QGP}{1+\eta g}=L,\nonumber\\
  \left[\Id-\eta PA^{-1}G\right]Q
  &=Q-\frac{\eta PGQ}{1+\eta g}=L^\dagger
\end{align}
give
\begin{equation}
  D\Phi[X]=LXL^\dagger,
\end{equation}
which proves Eq.~\eqref{eq:event-tangent}.  The derivation assumes
$1+\eta g>0$, exactly the condition that the selected branch has nonzero Born
probability.

\section{Monte Carlo form of the Born identity}
\label{app:monte-carlo}

For $S$ independent baseline trajectories $\bxi_a\sim p_0$, the direct
unbiased estimator is
\begin{equation}
  \widehat\chi^{\Born}_{\rm direct}
  =\frac1S\sum_{a=1}^S
  \left[\chi_{\bxi_a}^{\cond}
  +O_{\bxi_a}S_{\bxi_a}\right].
  \label{eq:mc-estimator}
\end{equation}
Because $\E[S]=0$, an exactly unbiased, sample-centered form uses the ordinary
sample covariance,
\begin{equation}
  \widehat\chi^{\Born}_{\rm cov}
  =\overline{\chi^{\cond}}
   +\frac{1}{S-1}\sum_{a=1}^S
    (O_{\bxi_a}-\bar O)(S_{\bxi_a}-\bar S).
  \label{eq:mc-covariance-estimator}
\end{equation}
Centering only $O$ but retaining the prefactor $1/S$ introduces an
$O(1/S)$ bias; Eq.~\eqref{eq:mc-covariance-estimator}, or an equivalent
leave-one-out construction, removes it.  The direct estimator
Eq.~\eqref{eq:mc-estimator} can have larger variance because a finite sample
generally has $\bar S\ne0$.

Sources and late origins inside trajectory $a$ should first be combined into a
single vector-valued observation $Y_a$.  Bootstrap resampling acts on the
$S$ complete $Y_a$ blocks, preserving all within-trajectory correlations and
the paired two-wall structure.

\section{Checklist for reading an H2 archive}
\label{app:archive-checklist}

Before interpreting a saved response product, verify:
\begin{enumerate}
  \item engine hash, covariance convention, source-kick sign, and schema;
  \item physical geometry, explicit wall bonds, wall windows, and orientation
  labels;
  \item independent trajectory count and shard/sample indices;
  \item initialization, total duration, origin cycles, and response horizon;
  \item random schedule law and whether schedules/outcomes were frozen in the
  derivative;
  \item derivative method: finite difference or exact tangent;
  \item finite-difference $\epsilon$ audit, if applicable;
  \item conditional, score, and total-response arrays separately;
  \item minimum Born denominator, replay residual, and censored branch counts;
  \item wall retention, bulk leakage, and total-charge drift;
  \item native Fourier resolution, time window, zero padding, and wedge
  definition;
  \item whole-trajectory bootstrap and paired-wall reversal test.
\end{enumerate}

\bibliographystyle{apsrev4-2}
\bibliography{references}

\end{document}
